\documentclass[11pt,a4paper]{article}
\usepackage[margin=1in]{geometry}
\usepackage{amsmath,amssymb,amsthm}
\usepackage{algorithm}
\usepackage{algpseudocode}
\usepackage{booktabs}
\usepackage{hyperref}

\theoremstyle{definition}
\newtheorem{definition}{\textbf{Definition}}[section]
\newtheorem{theorem}{\textbf{Theorem}}[section]
\newtheorem{lemma}[theorem]{\textbf{Lemma}}
\newtheorem{remark}{\textbf{Remark}}[section]

\title{\textbf{Knowledge Distillation via Soft Targets and Dark Knowledge Transfer}}
\author{\textbf{Team Scaibu} \\ \small Email: \texttt{jaydeep.9664920749@gmail.com} \\ \small \textbf{Architecture and Core Engine Team}}
\date{\textbf{October 2026}}

\begin{document}
\maketitle

\section{Definitions and Cognitive Mental Models}

\subsection{Formal Mathematical Definition}
\textbf{Definition (Knowledge Distillation via Soft Temperature Targets):}
Knowledge distillation (formulated by Hinton, Vinyals, \& Dean, 2015) transfers representational capacity and inductive bias from a large, high-capacity teacher network $f_{\text{teacher}}$ to a compact, low-latency student network $f_{\text{student}}$. Rather than training solely on one-hot hard labels $y \in \{0, 1\}^C$, the student matches the softened predictive probability distribution of the teacher across all classes.

Given unnormalized prediction logits $z^t \in \mathbb{R}^C$ from the teacher and $z^s \in \mathbb{R}^C$ from the student, temperature scaling by parameter $T > 1$ softens both output distributions:
{\boldmath
\begin{equation}
p_i^t(T) = \frac{\exp(z_i^t / T)}{\sum_{j=1}^C \exp(z_j^t / T)}, \quad p_i^s(T) = \frac{\exp(z_i^s / T)}{\sum_{j=1}^C \exp(z_j^s / T)}
\end{equation}
}
The total distillation loss is a convex combination of the Kullback-Leibler (KL) divergence between softened distributions and standard cross-entropy loss against ground truth:
{\boldmath
\begin{equation}
\mathcal{L}_{\text{KD}} = \alpha T^2 \cdot D_{\text{KL}}\left( p^t(T) \,\|\, p^s(T) \right) + (1 - \alpha) \cdot \mathcal{L}_{\text{CE}}\left( y, \, p^s(1) \right)
\end{equation}
}
where the scaling factor $T^2$ compensates for the fact that the gradients of the soft target loss scale as $1/T^2$ when temperature $T$ increases, ensuring balanced gradient magnitudes between the distillation and classification losses.

\subsection{Expanded Non-Technical Definition (Intuition for Anyone)}
\begin{enumerate}
    \item \textbf{Dark Knowledge Revealed}: Consider a picture of a BMW car. The hard label simply says ``Car''. A human expert or large teacher model knows: it is 90\% Car, but also 8\% Truck, 1.9\% Motorcycle, and 0.0001\% Carrot! Those relative non-target probabilities (``dark knowledge'') explain how concepts in reality relate to each other.
    \item \textbf{A Master-Apprentice Relationship}: Rather than forcing the student to memorize the answer key to a test, the student watches how the master teacher reasons through subtle confusions.
    \item \textbf{10x Serving Acceleration}: A 500-million parameter student trained via distillation frequently matches or exceeds 95\% of a 70-billion teacher's accuracy, cutting inference latency by an order of magnitude.
    \item \textbf{Temperature as a Magnifying Glass}: At standard temperature $T=1$, the winning class dominates with 99.9\% probability, hiding the dark knowledge. Raising temperature to $T=3$ flattens the distribution, bringing the subtle nuances into clear view for the student.
\end{enumerate}

\subsection{Child-Friendly Mental Model (Physical Metaphor)}
Imagine an Olympic chef teaching an apprentice:
\begin{itemize}
    \item \textbf{Hard Labels}: Handing the student a plate and saying ``This is soup.'' (True or False).
    \item \textbf{Soft Targets}: Letting the student taste the soup with an ultra-sensitive spoon ($T=2$), revealing: ``Notice that hint of nutmeg? Notice that subtle touch of ginger?'' The student learns all the secret ingredients that made the soup delicious!
\end{itemize}

\section{Core Mathematical Equations and Definitions}

\subsection{Distillation Loss Formulation}
{\boldmath
\begin{equation}
\mathcal{L}_{\text{KD}} = \alpha T^2 \sum_{i=1}^C p_i^t(T) \log \left(\frac{p_i^t(T)}{p_i^s(T)}\right) + (1 - \alpha) \left( - \sum_{i=1}^C y_i \log p_i^s(1) \right)
\end{equation}
}
\begin{itemize}
    \item \textbf{Formal Definition}: Multi-objective loss coupling scaled relative entropy with hard cross-entropy.
    \item \textbf{What It Means}: Guides the student to simultaneously predict the correct class and imitate teacher similarity structures.
    \item \textbf{Why It Is Important}: Foundational optimization objective for model compression and student distillation.
\end{itemize}

\subsection{Gradient Scaling Factor $T^2$}
{\boldmath
\begin{equation}
\frac{\partial \mathcal{L}_{\text{soft}}}{\partial z_i^s} \approx \frac{1}{T} (p_i^s(T) - p_i^t(T)) \xrightarrow{T \to \infty} \frac{1}{C T^2} (z_i^s - z_i^t)
\end{equation}
}
\begin{itemize}
    \item \textbf{Formal Definition}: Asymptotic gradient of the soft-target cross-entropy w.r.t. student logits.
    \item \textbf{What It Means}: Explains why multiplying by $T^2$ is necessary to keep gradient magnitudes commensurate with hard cross-entropy.
    \item \textbf{Why It Is Important}: Prevents catastrophic vanishing of distillation gradients during high-temperature training.
\end{itemize}

\section{Rigorous Mathematical Analysis Checklist}

\begin{itemize}
    \item \textbf{Notation}: $z^s, z^t \in \mathbb{R}^C$ (logits), $p^s(T), p^t(T)$ (probabilities at temperature $T$), $y \in \{0, 1\}^C$ (ground truth), $\alpha \in [0, 1]$, $T > 0$.
    \item \textbf{Epistemic Status}: Hinton et al. (NeurIPS 2015). Seminal foundational technique across neural compression.
    \item \textbf{Dependency Graph}: Student \& Teacher Logits $\to \text{Divide by } T \to \text{Softmax} \to \text{KL Divergence} \to \text{Multiply by } T^2 \to \text{Add } (1-\alpha)\text{CE}$.
    \item \textbf{Dimensions}: Logit dimension $C$ (class vocabulary / label size).
    \item \textbf{Regimes}: Edge deployment, real-time routing models, student distilled small language models (e.g. DistilBERT, Gemma-2-2B).
    \item \textbf{Invariants}: When $p^s(T) \equiv p^t(T)$, the distillation KL divergence is identically $0$.
    \item \textbf{Isomorphisms}: At $T \to \infty$, distillation is isomorphic to mean squared error (MSE) matching on zero-mean logits.
\end{itemize}

\section{Theoretical Theorems and Formal Proofs}

\begin{theorem}[High-Temperature Equivalence to Logit Matching]
Let $z^t, z^s \in \mathbb{R}^C$ be zero-mean teacher and student logits ($\sum_i z_i^t = \sum_i z_i^s = 0$). In the limit of infinite temperature $T \to \infty$, the gradient of the scaled soft-target distillation loss $\mathcal{L}_{\text{soft}} = T^2 D_{\text{KL}}(p^t(T) \| p^s(T))$ with respect to student logit $z_i^s$ converges to the gradient of Mean Squared Error:
{\boldmath
\begin{equation}
\lim_{T \to \infty} \frac{\partial \mathcal{L}_{\text{soft}}}{\partial z_i^s} = \frac{1}{C} (z_i^s - z_i^t) = \frac{\partial}{\partial z_i^s} \left( \frac{1}{2C} \|z^s - z^t\|_2^2 \right)
\end{equation}
}
\end{theorem}

\begin{proof}
Consider the Taylor expansion of the exponential function for large temperature $T$:
{\boldmath
\begin{equation}
\exp\left(\frac{z_i}{T}\right) = 1 + \frac{z_i}{T} + \frac{z_i^2}{2 T^2} + \mathcal{O}(T^{-3})
\end{equation}
}
Summing over all classes $C$, since $\sum_{j=1}^C z_j = 0$:
{\boldmath
\begin{equation}
\sum_{j=1}^C \exp\left(\frac{z_j}{T}\right) = C + \frac{1}{T} \sum_{j=1}^C z_j + \mathcal{O}(T^{-2}) = C + \mathcal{O}(T^{-2})
\end{equation}
}
Thus, the softmax probability simplifies to:
{\boldmath
\begin{equation}
p_i(T) = \frac{1 + z_i / T}{C} + \mathcal{O}(T^{-2})
\end{equation}
}
The gradient of $D_{\text{KL}}(p^t \| p^s)$ with respect to student logit $z_i^s$ is given by:
{\boldmath
\begin{equation}
\frac{\partial}{\partial z_i^s} D_{\text{KL}}(p^t \| p^s) = \frac{1}{T} (p_i^s(T) - p_i^t(T))
\end{equation}
}
Substituting the first-order approximations of $p_i^s(T)$ and $p_i^t(T)$:
{\boldmath
\begin{equation}
p_i^s(T) - p_i^t(T) = \frac{1 + z_i^s / T}{C} - \frac{1 + z_i^t / T}{C} + \mathcal{O}(T^{-2}) = \frac{z_i^s - z_i^t}{C T} + \mathcal{O}(T^{-2})
\end{equation}
}
Multiplying by the scale factor $T^2$:
{\boldmath
\begin{equation}
\frac{\partial \mathcal{L}_{\text{soft}}}{\partial z_i^s} = T^2 \left( \frac{1}{T} (p_i^s(T) - p_i^t(T)) \right) = T \left( \frac{z_i^s - z_i^t}{C T} + \mathcal{O}(T^{-2}) \right) = \frac{z_i^s - z_i^t}{C} + \mathcal{O}(T^{-1})
\end{equation}
}
Taking the limit as $T \to \infty$:
{\boldmath
\begin{equation}
\lim_{T \to \infty} \frac{\partial \mathcal{L}_{\text{soft}}}{\partial z_i^s} = \frac{1}{C} (z_i^s - z_i^t)
\end{equation}
}
This matches the exact gradient of $\frac{1}{2C} \sum_{j=1}^C (z_j^s - z_j^t)^2$, proving that at high temperatures, soft target distillation is mathematically equivalent to penalized logit regression.
\end{proof}

\begin{theorem}[Non-Negativity and Convexity of Knowledge Distillation Loss]
For any fixed teacher probability vector $p^t \in \Delta^C$ and student distribution $p^s \in \Delta^C$, the scaled distillation divergence $\mathcal{L}_{\text{distill}} = T^2 D_{\text{KL}}(p^t \| p^s)$ satisfies:
{\boldmath
\begin{equation}
\mathcal{L}_{\text{distill}} \ge 0
\end{equation}
}
with equality if and only if $p^s = p^t$ almost everywhere, and is strictly convex in $p^s$.
\end{theorem}

\begin{proof}
By Gibbs' inequality and Jensen's inequality applied to the negative logarithm:
{\boldmath
\begin{equation}
- D_{\text{KL}}(p^t \| p^s) = \sum_{i=1}^C p_i^t \log \left(\frac{p_i^s}{p_i^t}\right) \le \log \left( \sum_{i=1}^C p_i^t \frac{p_i^s}{p_i^t} \right) = \log\left( \sum_{i=1}^C p_i^s \right) = \log(1) = 0
\end{equation}
}
Multiplying by $-1$ establishes $D_{\text{KL}}(p^t \| p^s) \ge 0$.
Since $T > 0$, $T^2 > 0$, and thus $\mathcal{L}_{\text{distill}} = T^2 D_{\text{KL}}(p^t \| p^s) \ge 0$.
Equality holds if and only if the argument $\frac{p_i^s}{p_i^t} = 1$ for all $i$ with $p_i^t > 0$, i.e. $p^s \equiv p^t$.
Strict convexity follows from the strict convexity of $-\log(x)$ on $\mathbb{R}_{++}$, guaranteeing a unique global minimum when the student perfectly mimics the teacher.
\end{proof}

\section{Foundational Architectural Questions and Epistemic Meta-Questions}

\subsection{Architectural Question 1: Optimal Temperature Tuning}
\textbf{Question:} How should temperature $T$ be chosen in practical distillation setups?\\
\textbf{Answer:} Temperature typically ranges between $T = 2.0$ and $T = 5.0$. If the student is very small relative to the teacher, low temperatures ($T \approx 1.5\text{--}2.0$) are preferable so the student focuses on prominent high-confidence classes. If the student has substantial capacity, higher temperatures ($T \approx 4.0\text{--}8.0$) expose fine-grained dark knowledge structures without confusing the student.

\textbf{Meta-Question:} Should intermediate hidden states also be distilled?\\
\textbf{Answer:} Yes. In \textbf{Feature-Based Distillation} (FitNets), mean squared error is computed between student intermediate activations $W_{\text{proj}} h^s$ and teacher activations $h^t$, guiding the student throughout its depth rather than solely at the output logits.

\subsection{Architectural Question 2: Online Distillation and Co-Training}
\textbf{Question:} Can distillation work when a large teacher model does not exist beforehand?\\
\textbf{Answer:} Yes. In \textbf{Deep Mutual Learning} (online distillation), an ensemble of multiple peer student models train concurrently from scratch, each treating the other peers as teachers by minimizing bidirectional KL divergence, yielding models that generalize significantly better than isolated models.

\textbf{Meta-Question:} How does knowledge distillation interact with model quantization?\\
\textbf{Answer:} Distillation is exceptionally effective during Quantization-Aware Training (QAT). The full-precision FP16 teacher acts as an anchor, preventing the 4-bit or 8-bit student from suffering quantization noise degradation.

\section{Algorithmic Implementation Specification}

\begin{algorithm}[H]
\caption{Soft-Target Knowledge Distillation Loss Evaluation}
\begin{algorithmic}[1]
\Require Student logits $\mathbf{z}^s \in \mathbb{R}^C$, teacher logits $\mathbf{z}^t \in \mathbb{R}^C$, target label $y \in \{0, \dots, C-1\}$, temperature $T$, weight $\alpha$
\Ensure Combined loss $\mathcal{L}$, KL loss $\mathcal{L}_{\text{distill}}$, CE loss $\mathcal{L}_{\text{hard}}$, probabilities $\mathbf{p}^s(T), \mathbf{p}^t(T)$
\State Compute softened probabilities:
\State $\mathbf{p}^s(T) \gets \operatorname{Softmax}(\mathbf{z}^s / T), \quad \mathbf{p}^t(T) \gets \operatorname{Softmax}(\mathbf{z}^t / T)$
\State Initialize $\text{kl} \gets 0.0$
\For{$i = 0$ \textbf{to} $C-1$}
    \State $t_i \gets \max(10^{-12}, \mathbf{p}^t_i(T)), \quad s_i \gets \max(10^{-12}, \mathbf{p}^s_i(T))$
    \State $\text{kl} \gets \text{kl} + t_i \cdot \ln(t_i / s_i)$
\EndFor
\State $\mathcal{L}_{\text{distill}} \gets T^2 \cdot \max(0.0, \text{kl})$
\State Compute hard student distribution: $\mathbf{p}^s(1) \gets \operatorname{Softmax}(\mathbf{z}^s)$
\State $\mathcal{L}_{\text{hard}} \gets -\ln(\max(10^{-12}, \mathbf{p}^s_y(1)))$
\State $\mathcal{L} \gets \alpha \cdot \mathcal{L}_{\text{distill}} + (1 - \alpha) \cdot \mathcal{L}_{\text{hard}}$
\State \Return $(\mathcal{L}, \mathcal{L}_{\text{distill}}, \mathcal{L}_{\text{hard}}, \mathbf{p}^s(T), \mathbf{p}^t(T))$
\end{algorithmic}
\end{algorithm}

\section{Big-$\Theta$ Complexity Analysis}

\begin{itemize}
    \item \textbf{Softmax Evaluation Time}: $\Theta(C)$ operations for student and teacher at temperature $T$.
    \item \textbf{KL Divergence Evaluation Time}: $\Theta(C)$ operations.
    \item \textbf{Hard Cross-Entropy Time}: $\Theta(C)$ for $T=1$ normalization.
    \item \textbf{Total Time Complexity}: $\Theta(C)$ operations per training sample.
    \item \textbf{Memory Storage Overhead}: $\Theta(C)$ floats for probability distributions.
\end{itemize}

\section{Single Canonical Repository Reference}

The complete deterministic scalar implementation, verification suite, and unit test benchmarks are published at:
\begin{center}
\url{https://github.com/Chief-Strategist-J/llm-obs-infra}
\end{center}

\end{document}
